%%%%%%%%%%%%%%%%%%%%%%%%%%%%%%%%%%%%%%%%%
% Edwin Mei's CV
%
% Important note:
% This template requires the resume.cls file to be in the same directory as the
% .tex file. The resume.cls file provides the resume style used for structuring the
% document.
%
%%%%%%%%%%%%%%%%%%%%%%%%%%%%%%%%%%%%%%%%%

%----------------------------------------------------------------------------------------
%	PACKAGES AND OTHER DOCUMENT CONFIGURATIONS
%------------------------------------------------\----------------------------------------

\documentclass{resume} % Use the custom resume.cls style
\usepackage{xeCJK}
% \setCJKmainfont{FandolSong}
\setCJKmainfont[BoldFont={Microsoft YaHei UI Bold}]{楷体}
\usepackage[left=0.75in,top=0.6in,right=0.75in,bottom=0.6in]{geometry} % Document margins

\name{梅越} % Name
\address{主页：https://tuoye25.github.io/} % Home
\address{名片：https://bonjour.bio/edwinm/} % Busines Card
\address{电话：(+86)~$\cdot$~18964070908} % Phone number
\address{邮箱：edwin\_mei-yue@foxmail.com} % Email

\begin{document}

\begin{rSection}{个人简介}
  这里是梅越，目前担任高级AI开发工程师，同时以CTO身份参与创业项目的技术决策与架构设计。我的研究方向聚焦于深度学习与模型优化，在知识蒸馏与计算机视觉领域有深入探索，并以第一作者在 ICICML 2023 发表相关论文。工业经验方面，我曾任职于微软（中国）Azure 云架构团队，并在富谙智科技主导数据分析与全栈开发，具备云原生架构、AI 工程化落地及跨团队协作的完整能力。技术栈覆盖 PyTorch、TensorFlow、LangGraph 等主流框架，以及 Claude、DeepSeek、Multica 等 AI Agent 工具链。期待将学术积累与工程实践结合，在 AI 软件研发的岗位上持续创造价值。
\end{rSection}

\begin{rSection}{学历}

  {\bf Roux研究院，美国东北大学} \hfill {\em 06.2025} \\
  计算机科学硕士 \hfill GPA: \textbf {4.0}

  {\bf 俄亥俄州立大学，美国} \hfill {\em 06.2019} \\
  电气与计算机工程学士 \hfill 数学辅修

\end{rSection}

\begin{rSection}{工作}

  \begin{rSubsection}{群贤集团}{06.2026 - 至今}{高级AI开发工程师}{上海}
    \item 主导公司数字化转型战略，负责AI技术架构设计与落地实施
    \item 桂林工厂AI数据底座：设计并构建工业级AI数据底座，整合多源异构生产数据，负责模型训练与性能优化，实现生产流程智能化升级
    \item 公司ERP系统设计与优化：主导ERP系统智能化改造，将AI能力嵌入核心业务流程，提升供应链管理与资源配置效率
    \item OA平台AI化转型：推动公司OA系统智能化升级，集成智能审批、知识问答、流程自动化等AI功能，显著提升办公效率
  \end{rSubsection}

  \begin{rSubsection}{数字工厂（副业）}{01.2026 - 至今}{联合创始人 \& CTO}{上海}
    \item 主导构建基于多智能体协作的分布式 AI 工厂系统，负责整体技术架构设计与 AI 模型编排
    \item 设计并落地混合云-本地推理架构，整合 Claude Code、DeepSeek V4 Pro、Qwen（蒸馏 Opus 4.7）等多型号大模型，通过 Ollama 实现本地低成本推理，优化响应速度与算力成本
    \item 基于 Multica + OpenClaw 搭建多 Agent 协作框架，调度 Hermes、QoderWake 等数字员工并行处理复杂任务，实现游戏设计业务的自动化流水线
    \item 从 0 到 1 组建技术体系，覆盖 AI 推理、任务编排、多模型路由等核心模块
    \item 业务从游戏设计拓展至团购等多条线，系统具备良好的可迁移性与横向扩展能力
  \end{rSubsection}

  \begin{rSubsection}{富谙智信息技术有限公司}{12.2021 - 07.2022}{数据分析研发师}{上海}
    \item 在敏捷开发环境中以全栈工程师的角色负责研发与改进核心库 FNZOneX
    \item 为顾问组设计新的用户组结构并修改由 Alchemist 基于用户权限矩阵生成的定义脚本
    \item 为不同开发/测试环境制作运行 SQL 脚本组
    \item 在浏览器中通过检查的功能（元素，主控，资源，网络）调试 JavaScript 报错；通过 VS 调试模式定位
    aspx/aspx.vb 报错
    \item 设置，修改和分享 IIS 服务器配置用于开发和演示
    \item 通过和国际 IT 部门以及跨团队部门合作以解决管辖相关问题
    \item 解决由测试员发现的显示和发布问题， 并且和测试员一起功能审查前自测并制订测试流程
  \end{rSubsection}

  \begin{rSubsection}{上海微创软件股份有限公司（微软大中华区）}{10.2019 - 12.2021}{云架构工程师 \& 技术顾问}{上海}
    \item 专注于微软云 Azure 中四个产品：Monitor, Log Analytics, Automation and Application Insights 的技术支持
    \item 定制警报配置、查询搜索、数据采样和多样性能计数器采集
    \item 捕获并分析 NetMon、ProcMon 和 Fiddler 跟踪，以解决代理、混合工作程序和门户问题；通过 Azure
    仪表板、度量资源管理器和工作簿可视化数据
    \item 使用 Visual Studio/VS 代码在.Net/.Net Core/Java/Python/JavaScript 中调试 Application Insights SDK 配置脚本
    \item 帮助设计、修改和调试 PowerShell/Python/PS 工作流/图形运行手册和 DSC 配置；提供有关合作伙伴集成和工作区解决方案的建议，如更新管理、自动启动/停止虚拟机、更改跟踪和资源清册
    \item 对接 Adobe，Intel 和 Amazon 等关键客户并专门处理风险或长期案例升级
    \item 负责保护上海，无锡和北京在内的 12 人团队并解决跨团队争议
    \item 与领队以及协作技术顾问合作进行案例分析，技术面试和团队规划
  \end{rSubsection}

\end{rSection}

\begin{rSection}{项目}

  \begin{rSubsection}{银河卷积神经网络的数据集蒸馏}{09.2024 - 12.2025}{导师：Tala Talaei Khoei}{NEU}
    \item 通过使用最新 Galaxy Zoo 数据集的提取和整理版本，在星系形态分类等任务中展现出卓越的性能
    \item 仔细调整统计显著性阈值以确定验证损失趋势，以优化性能
    \item 使用 STM 来防止过拟合并加速学习过程
    \item 提供更高的分类准确性和更高效的训练，同时减少类别失衡并提高整体数据质量
  \end{rSubsection}

  \begin{rSubsection}{深度强化学习的知识蒸馏}{07.2023 - 09.2023}{导师：Mark Vogelsberger}{MIT}
    \item 通过 pytorch 提供的深度 Q 网络搭建深度强化学习的教师与学生模型
    \item 在 Gymnasium 提供的 CartPole 环境中进行实验比较
    \item 通过改变温度和回归次数等参数来测试知识蒸馏的效率
    \item 在 CartPole, MountainCar, Acrobot, LunarLander and Pendulum 环境中比较不同的马可夫决策过程
    \item 在更复杂的深度回归 Q 网络中实现知识蒸馏
    \item 本项目的研究成果已形成学术论文，发表于机器学习顶级会议 ICICML 2023：\\https://doi.org/10.1109/ICICML60161.2023.10424889
  \end{rSubsection}

  \begin{rSubsection}{毕业设计：太阳能电池板检测仪}{01.2019 - 04.2019}{导师：Haskell Jac Fought}{OSU}
    \item 设计并修改软件部分的需求，特征和行为流程图
    \item 在 Android Studio 上设计，编程并测试蓝牙串行通讯模块
    \item 从 Arduino 机械码转换已有代码至安卓
  \end{rSubsection}

  \begin{rSubsection}{AEV 设计，搭建与调试}{01.2016-05.2016}{导师：Lauren Corrigan}{OSU}
    \item 通过 Arduino 集成开发编译器编程
    \item 用 Solid Works 设计 Arduino 的部件，精准控制其尺寸和形状
    \item 测试并记录 Arduino 的表现，控制变量以提升其空气动力效率
  \end{rSubsection}

\end{rSection}

\begin{rSection}{实习 \& 兼职}

  \begin{rSubsection}{Roux研究院}{01.2025 - 05.2025}{机器学习助教}{美国}
    \item 辅导学生了解不同类型的的机器学习模型以及相关技巧
    \item 与教授讨论课程分配；在办公时间回答问题并检阅项目
  \end{rSubsection}

  \begin{rSubsection}{俄州大电气与计算机系实验室}{09.2018 - 12.2018}{实验助教 \& 阅卷人}{美国}
    \item 实验期间指导学生练习 TI MPS430 系列微控制器编码
    \item 批改每周测试和实验报告
  \end{rSubsection}

  \begin{rSubsection}{华虹计通}{05.2018 - 07.2018}{软件研发工程师}{上海}
    \item 学习自动售检票系统系统构架，数据流设计，运作原则以及编码规则
    \item 改进西安地铁 1 至 4 号线地铁 UI 交互界面并更新最新地铁线路数据
    \item 从甲骨文 SQL 至 DB2 的数据库转移
  \end{rSubsection}

  \begin{rSubsection}{华信能源}{05.2017 - 07.2017}{信息系统工程师}{上海}
    \item 改善公司前端 OA 系统，新增人性化时间表并与公司 app (Android / IOS)连接
    \item 设计 Access VBA 宏指令集，自动计算各部门电话账单并进行简单的筛选、报表统计和导出
  \end{rSubsection}

\end{rSection}

\begin{rSection}{技能}

  \begin{tabular}{ @{} >{\bfseries}l @{\hspace{6ex}} p{14cm} }
    编程语言       & Python, C\#, Java, C++, Golang, Shell, JavaScript, M, VBA, VHDL \\
    机器学习框架     & PyTorch, TensorFlow, Keras, Scikit-learn, LangGraph             \\
    模型与AI工具    & Claude, Codex, DeepSeek, Pi, Wan, Flux, Kling, Seedance         \\
    AI Agent框架 & Multica, HiClaw, PaperClips, Hermes                             \\
    模型优化技术     & 知识蒸馏 (Knowledge Distillation), 模型量化, 剪枝, 推理加速                   \\
    数据库        & Microsoft SQL, KQL, MySQL, MongoDB, SQLite                      \\
    云平台与架构     & Azure, AWS, 分布式系统, 微服务                                          \\
    开发工具       & VS, VS Code, Git, Docker, SQL Server, Kusto, TailScale          \\
    硬件基础       & Arduino, Multiplexers, PLDs, Clocked Sequential Circuits, VLSI  \\
  \end{tabular}

\end{rSection}

\begin{rSection}{课程}

  \begin{tabular}{ @{} >{\bfseries}l @{\hspace{6ex}} p{12cm} }
    计算机 & 机器学习, 数据库管理系统, 算法, 云计算, 软件开发与设计, 操作系统, 计算机体系结构与设计, 计算机组成与嵌入式系统          \\
    电气  & 模拟系统与电路, 信号与系统, 数字逻辑, 微控制器, 半导体电子器件, 光学工程, 反馈控制系统, 高级数字设计, VLSI数字信号处理系统 \\
    数学  & 科学家与工程师微积分, 线性代数, 高等数学, 数值分析, 偏微分方程                                     \\
    其他  & 工程经济决策, 工程伦理, 概率论与数理统计                                                  \\
  \end{tabular}

\end{rSection}

\begin{rSection}{兴趣爱好}
  工作之外，我是个动静皆宜的人。动的方面，我热爱匹克球和羽毛球，享受球场上快节奏对抗和团队配合的默契；静的方面，我是密室逃脱、剧本杀和桌游的资深爱好者，喜欢在沉浸式剧情和烧脑推理中切换思维、体验不同视角。这些爱好让我深信：好的 AI 产品也像一场精心设计的游戏——需要对用户心理的深度共情，以及对细节体验的极致打磨。期待未来团建时和大家来一场酣畅淋漓的切磋。
\end{rSection}

\end{document}